\subsubsection{Hàm chi phí không đồng nhất}
\label{sec:bloxorz-action-cost}
Gọi $O(s)$ là tập ô bị chiếm trong trạng thái $s$ và $F$ là tập ô dễ vỡ. Chi phí
của cạnh hợp lệ $(s,a,s')$ được định nghĩa:
\begin{equation}
  c(s,a,s')=1+\left|\bigl(O(s')\setminus O(s)\bigr)\cap F\right|.
  \label{eq:bloxorz-action-cost}
\end{equation}
Mỗi lần lăn có chi phí nền $1$; mỗi ô dễ vỡ mới tiếp xúc cộng thêm $1$.
Khối chiếm nhiều nhất hai ô nên chi phí thuộc $\{1,2,3\}$ và
$c_{\min}=1>0$. Chính sách nằm tại \path{bloxorz.core.costs};
\texttt{Board.step\_cost} chỉ gọi chính sách này. Điều chỉnh trọng số không đòi
hỏi sửa luật chuyển trạng thái hoặc vòng lặp tìm kiếm.

Phụ phí thể hiện ưu tiên hạn chế tiếp xúc bề mặt dễ vỡ, theo ví dụ của đề bài.
Đây là mục tiêu tối ưu được thiết kế, không phải xác suất rơi: trò chơi vẫn xác
định. Hệ số $1$ tạo đánh đổi dễ giải thích, trong đó tránh một tiếp xúc tương
đương giảm một lần lăn thường; không lựa chọn hệ số bằng cách tối ưu kết quả
thực nghiệm. Lập phương đứng yên trên ô dễ vỡ không làm tăng chi phí mọi bước
của lập phương còn lại. Đích dịch chuyển được xét như tiếp xúc mới; cầu và công
tắc không có phụ phí riêng.

Hợp đồng hành động hiện hành không có cạnh đổi điều khiển độc lập. Vì vậy,
không thêm cạnh chi phí $0$ hoặc $1$ cho việc chọn lập phương; mọi cạnh di
chuyển thực sự vẫn có chi phí ít nhất $c_{\min}$. Tổng chi phí $g$ là tổng trọng số cạnh, còn độ
dài lời giải là \texttt{len(path)}; hai đại lượng không được đánh đồng.


Đối với hàm ước lượng dựa trên khoảng cách, cận số bước phải được nhân với
$c_{\min}$ để thành cận chi phí. Chẳng hạn, trên màn không dịch chuyển, nếu
$d$ giảm nhiều nhất hai ô sau mỗi lần lăn thì $h=c_{\min}d/2$ không vượt
chi phí còn lại. Khi có dịch chuyển, lập luận hình học này cần được thay thế
bằng một cận đã chứng minh; $h=0$ luôn là lựa chọn an toàn. Phần A* sẽ trình
bày thiết kế và kiểm chứng cụ thể.
